% !TeX root = ../main.tex
\section{引言}

电动汽车的长途出行需要沿高速公路获得连续、可靠的补能服务。换电站通过交付已经充好的电池缩短车辆的现场补能时间，其服务能力同时取决于车辆到站分布和站内电池库存。长途车辆可能在多个站点换电，前一次选站会改变后续可达站点；每次换电还会向站内退回一块具有特定 SOC 的电池，使车辆的换电安排与各站未来的充电需求相互影响。

运营中同时存在预约用户和随机到站用户。预约信息为需求安排提供依据，车辆实际进入高速公路后的位置、SOC 和预计到站时间又会不断更新。随机用户的到站时间和退回电池状态具有不确定性。与此同时，分时电价、设备功率和站级供电能力共同影响充电安排。因此，运营者需要结合实时信息选择预约用户的换电站点，并协调各站充电功率，以平衡服务收入、购电费用、路径调整和预约履约。

换电路径与充电功率具有直接联系。将车辆安排到某一站点，会改变该站的电池需求、退回 SOC 和后续充电负荷；提高某一站点的充电功率，也可能使原本库存紧张的换电安排获得可行的服务时间。本文将两类决策纳入同一滚动时域优化模型，在每轮预测域内统一评价其运营影响。充电功率按较短周期更新，换电路径按较长周期更新，以适应站内状态变化并保持用户换电安排的稳定。

有限预测域会截断电池库存和预约服务的长期影响。预测结束时仍在充电的电池可能服务于后续需求，已经确定的预约行程也可能在预测域外继续换电。终端价值用于评价这些状态对后续运营的影响。已有研究通过近似价值迭代构造预测控制的终端成本，并分析其与预测长度及价值近似误差的关系\cite{moreno2022terminal}。本文结合高速公路换电场景，将终端库存与预测域外未完成服务共同纳入状态价值函数，通过实际运营轨迹训练，并在每轮优化中计算本轮决策对应的终端价值。

本文的研究内容包括三个方面。首先，利用 SOC 分档候选网络描述预约用户的可达换电路径，建立换电路径与充电功率的联合优化模型。其次，在模型中表达满电交付、退回电池充电、预约优先、等待窗口及路径调整费用，使决策与电池和用户状态的变化相互对应。最后，构造包含终端库存和未完成服务的价值函数，并通过联合优化消融、终端价值消融及敏感性分析评价各组成部分的作用。

图~\ref{fig:framework} 给出整体方法。第 2 节介绍问题与基本假设；第 3 节建立候选网络、日前换电路径和滚动时域优化模型；第 4 节介绍终端价值学习；第 5 节介绍数值实验；第 6 节总结全文。附录给出分段线性函数和相关变量乘积的等价约束。

\begin{figure}[!htbp]
\centering
\begin{tikzpicture}[
  box/.style={draw=blue!45!black,rounded corners=2pt,align=center,
    text width=3.8cm,minimum height=1.1cm,fill=blue!3,font=\small},
  arr/.style={-{Latex[length=2mm]},thick,draw=blue!55!black}]
  \node[box] (obs) {实时状态与预测\\车辆、请求、电池、电价};
  \node[box,right=0.75cm of obs] (opt) {滚动时域优化\\换电路径与充电功率};
  \node[box,right=0.75cm of opt] (act) {当前时间段执行\\服务、充电与状态更新};
  \node[box,below=1.1cm of opt] (value) {RL 终端价值\\终端库存与未完成服务};
  \draw[arr] (obs)--(opt);
  \draw[arr] (opt)--(act);
  \draw[arr] (value)--node[right,font=\small]{价值函数}(opt);
  \draw[arr] (act.south)|-node[pos=0.3,right,font=\small]{实际收益与轨迹}(value.east);
  \draw[arr] (act.north)--++(0,0.6)-|node[pos=0.5,above,font=\small]{下一时刻观测}(obs.north);
\end{tikzpicture}
\caption{换电路径与充电功率滚动时域优化方法}
\label{fig:framework}
\end{figure}

